\documentclass[11pt,a4paper]{article}

\usepackage[T1]{fontenc}
\usepackage[utf8]{inputenc}
\usepackage[spanish,es-noshorthands]{babel}
\usepackage{lmodern}
\usepackage{amsmath,amssymb,amsthm,mathtools}
\usepackage[margin=2.5cm]{geometry}
\usepackage{microtype}
\usepackage[colorlinks=true,linkcolor=black,citecolor=black,urlcolor=blue]{hyperref}

\newcommand{\dd}{\,\mathrm{d}}
\newcommand{\p}{\partial}
\newcommand{\sat}{\mathrm{sat}}
\newcommand{\lit}{\mathrm{lit}}
\newcommand{\ad}{\mathrm{ad}}
\newcommand{\Rv}{\mathcal{R}}

\theoremstyle{plain}
\newtheorem{proposicion}{Proposición}
\theoremstyle{definition}
\newtheorem{hipotesis}{Hipótesis}
\theoremstyle{remark}
\newtheorem*{obs}{Observación}

\title{\vspace{-1.4cm}
\textbf{Pendiente de presión en el tramo sin burbujas}\\[0.35em]
\large Promedio en la sección del momento 2D, con el divisor de compresibilidad}
\author{Ignacio Gómez Gómez}
\date{\today}

\begin{document}
\maketitle
\vspace{-0.7cm}

\begin{abstract}\noindent
En el tramo en que el agua sigue disuelta, el momento vertical del fundido,
promediado en la sección y evaluado sobre la parábola de Poiseuille, da
\[
\frac{\dd P}{\dd z}
=
-\,
\frac{\rho_m g+8\mu\bar u/R^2}
{1-\dfrac{4}{3}\,\bar u^2/c^2},
\qquad
c^2=\frac{K}{\rho_m}.
\]
El numerador es el peso más el roce de pared. El denominador es la inercia
que aparece al comprimir el fundido manteniendo el caudal de cada radio.
El factor \(4/3\) es \(\langle u^2\rangle/\bar u^2\) de esa parábola.
Un metro por encima de la saturación, el estado inicial de la marcha
diferencial queda determinado: \(\varphi\) es el volumen de Henry y
\(u_m(r)=u_g(r)\) conserva la parábola del caudal.
Esta nota escribe la cadena desde las ecuaciones reducidas hasta esas
expresiones.
\end{abstract}

\section{Ecuaciones de partida}

El conducto es el cilindro \(0\le r\le R\), \(H\le z\le 0\), con \(z\) hacia
arriba. En la reducción las velocidades radiales son nulas. Mientras el agua
sigue disuelta,
\begin{equation}
\varphi\equiv 0,
\qquad
\Gamma\equiv 0,
\qquad
F_{mg}\equiv 0,
\qquad
H\le z\le z_{\sat}.
\label{eq:phi0}
\end{equation}
La masa del fundido en cada radio es
\begin{equation}
\frac{\p}{\p z}\bigl(\rho_m u\bigr)=0.
\label{eq:masa}
\end{equation}
El momento vertical del fundido es
\begin{equation}
\rho_m u\frac{\p u}{\p z}
=
-\frac{\p P}{\p z}
-\rho_m g
+\frac{1}{r}\frac{\p}{\p r}\!\left(r\mu\frac{\p u}{\p r}\right)
+\frac{\p}{\p z}\!\left(\mu\frac{\p u}{\p z}\right).
\label{eq:mom}
\end{equation}
La componente radial del momento, con velocidad radial nula, da
\(\p P/\p r=0\). Por tanto
\begin{equation}
P=P(z).
\label{eq:Pdez}
\end{equation}
Cada término de la masa y del momento del gas lleva un factor \(\varphi\) o
\(F_{mg}\). Con~\eqref{eq:phi0} esas dos ecuaciones quedan en la identidad
\(0=0\), y \(u_g\) no entra en este tramo.

Las condiciones de borde del tramo son
\begin{equation}
\frac{\p u}{\p r}(0,z)=0,
\qquad
u(R,z)=0,
\qquad
P(H)=P_{\lit}+\Delta P.
\label{eq:bc}
\end{equation}

\begin{hipotesis}
\label{hip:cierre}
Se toman tres cierres, los mismos que producen la recta algebraica del modelo
1D.
\begin{enumerate}
\item \(\mu\) es constante en \(r\) y en \(z\).
\item El módulo de compresibilidad \(K\) es constante:
      \(\dd P=K\,\dd\rho_m/\rho_m\).
\item El perfil ya está desarrollado: tiene la forma parabólica que se
      construye en la sección~\ref{sec:parabola}, en todo
      \(H\le z\le z_{\sat}\).
\end{enumerate}
\end{hipotesis}

\section{Lo que fija la masa}

Como \(P=P(z)\) y \(\rho_m=\rho_m(P)\), la densidad depende solo de \(z\).
Integrar~\eqref{eq:masa} en \(z\) deja un caudal másico por unidad de área,
distinto en cada radio y constante a lo largo del conducto:
\begin{equation}
\rho_m(z)\,u(r,z)=q(r).
\label{eq:q}
\end{equation}
Derivando~\eqref{eq:q},
\begin{equation}
\frac{\p u}{\p z}
=
-\frac{u}{\rho_m}\frac{\dd\rho_m}{\dd z}.
\label{eq:du}
\end{equation}

\section{Compresibilidad}

Con la hipótesis~\ref{hip:cierre}.2,
\begin{equation}
\frac{\dd\rho_m}{\dd z}
=
\frac{\rho_m}{K}\frac{\dd P}{\dd z}.
\label{eq:eos}
\end{equation}
Se define \(c^2=K/\rho_m\), de modo que \(K=\rho_m c^2\).
Sustituyendo~\eqref{eq:eos} en~\eqref{eq:du},
\begin{equation}
\frac{\p u}{\p z}
=
-\frac{u}{K}\frac{\dd P}{\dd z}
=
-\frac{u}{\rho_m c^2}\frac{\dd P}{\dd z}.
\label{eq:duP}
\end{equation}
La inercia de~\eqref{eq:mom} queda
\begin{equation}
\rho_m u\frac{\p u}{\p z}
=
-\frac{u(r,z)^2}{c^2}\frac{\dd P}{\dd z}.
\label{eq:inercia}
\end{equation}
El divisor de la fórmula final sale de este término, al pasarlo al lado de
\(\dd P/\dd z\).

\section{Ecuación puntual}

Se lleva~\eqref{eq:inercia} a~\eqref{eq:mom}:
\begin{equation}
-\frac{u^2}{c^2}\frac{\dd P}{\dd z}
=
-\frac{\dd P}{\dd z}
-\rho_m g
+\frac{1}{r}\frac{\p}{\p r}\!\left(r\mu\frac{\p u}{\p r}\right)
+\frac{\p}{\p z}\!\left(\mu\frac{\p u}{\p z}\right).
\end{equation}
Al agrupar \(\dd P/\dd z\),
\begin{equation}
\frac{\dd P}{\dd z}
\left(1-\frac{u(r,z)^2}{c^2}\right)
=
-\rho_m g
+\frac{1}{r}\frac{\p}{\p r}\!\left(r\mu\frac{\p u}{\p r}\right)
+\frac{\p}{\p z}\!\left(\mu\frac{\p u}{\p z}\right).
\label{eq:puntual}
\end{equation}
El factor \(1-u^2/c^2\) depende de \(r\): en el eje \(u\) es máximo y en la
pared \(u=0\), así que ahí el factor vale \(1\). Una sola función \(P(z)\)
se obtiene cumpliendo~\eqref{eq:puntual} en promedio sobre la sección.

\section{Promedio en la sección}

Para una función \(f(r)\),
\begin{equation}
\langle f\rangle
=
\frac{2}{R^2}\int_0^R f(r)\,r\,\dd r.
\label{eq:media}
\end{equation}
El promedio de una cantidad que depende solo de \(z\) es ella misma.
Aplicando~\eqref{eq:media} a~\eqref{eq:puntual},
\begin{equation}
\begin{aligned}
\frac{\dd P}{\dd z}
\left(1-\frac{\langle u^2\rangle}{c^2}\right)
&=
-\rho_m g
+\left\langle
\frac{1}{r}\frac{\p}{\p r}\!\left(r\mu\frac{\p u}{\p r}\right)
\right\rangle
\\
&\qquad
+\left\langle
\frac{\p}{\p z}\!\left(\mu\frac{\p u}{\p z}\right)
\right\rangle.
\end{aligned}
\label{eq:prom}
\end{equation}

\section{El viscoseo radial, integrado}

Con \(\mu\) constante,
\begin{align}
\left\langle
\frac{1}{r}\frac{\p}{\p r}\!\left(r\mu\frac{\p u}{\p r}\right)
\right\rangle
&=
\frac{2}{R^2}\int_0^R
\frac{\p}{\p r}\!\left(r\mu\frac{\p u}{\p r}\right)\dd r
=
\frac{2}{R^2}
\left[
r\mu\frac{\p u}{\p r}
\right]_0^R.
\label{eq:Lr-int}
\end{align}
En el eje, \(r\,\p u/\p r\to 0\) por~\eqref{eq:bc}. En la pared queda
\begin{equation}
\left\langle L_r u\right\rangle
=
\frac{2\mu}{R}
\left.\frac{\p u}{\p r}\right|_{r=R}.
\label{eq:Lr-pared}
\end{equation}
La identidad~\eqref{eq:Lr-pared} vale para cualquier perfil que cumpla
\(u(R)=0\) y la simetría en el eje. El factor \(8\) aparece al evaluar la
derivada de la parábola.

\section{La parábola}
\label{sec:parabola}

Se busca un perfil \(u=A(z)+B(z)\,r^2\) que cumpla~\eqref{eq:bc}.
La simetría en \(r=0\) elimina un término lineal en \(r\).
La pared \(u(R,z)=0\) fija
\begin{equation}
u(r,z)=u_c(z)\left(1-\frac{r^2}{R^2}\right).
\label{eq:uc}
\end{equation}
La velocidad media es, con \(\eta=r/R\),
\begin{align}
\bar u(z)
&=
\langle u\rangle
=
\frac{2}{R^2}\int_0^R
u_c\left(1-\frac{r^2}{R^2}\right)r\,\dd r
=
2u_c\int_0^1(1-\eta^2)\eta\,\dd\eta
\nonumber\\
&=
2u_c\left[\frac{\eta^2}{2}-\frac{\eta^4}{4}\right]_0^1
=
2u_c\left(\frac12-\frac14\right)
=
\frac{u_c}{2}.
\label{eq:media-u}
\end{align}
Luego \(u_c=2\bar u\) y
\begin{equation}
u(r,z)=2\bar u(z)\left(1-\frac{r^2}{R^2}\right).
\label{eq:parabola}
\end{equation}
La derivada radial es
\begin{equation}
\frac{\p u}{\p r}
=
-\frac{4\bar u\, r}{R^2},
\qquad
\left.\frac{\p u}{\p r}\right|_{r=R}
=
-\frac{4\bar u}{R}.
\label{eq:dur}
\end{equation}
El esfuerzo en la pared es \(\tau_w=\mu\cdot 4\bar u/R\).
Sustituyendo~\eqref{eq:dur} en~\eqref{eq:Lr-pared},
\begin{equation}
\left\langle L_r u\right\rangle
=
\frac{2\mu}{R}\left(-\frac{4\bar u}{R}\right)
=
-\frac{8\mu\bar u}{R^2}.
\label{eq:F}
\end{equation}
El término \(8\mu\bar u/R^2\) es el roce \(\tau_w\cdot 2\pi R\) repartido en el
área \(\pi R^2\). En el balance 2D entra por~\eqref{eq:F}; no se suma un
\(F\) adicional.

Con~\eqref{eq:q} y~\eqref{eq:parabola}, el caudal medio
\(\bar q=\langle q\rangle\) cumple
\begin{equation}
\bar u(z)=\frac{\bar q}{\rho_m(z)}.
\label{eq:ubar}
\end{equation}
La forma parabólica se conserva en todo el tramo. La amplitud cambia solo
cuando cambia \(\rho_m\).

\section{El promedio de \texorpdfstring{\(u^2\)}{u al cuadrado}}

En el divisor de~\eqref{eq:prom} entra \(\langle u^2\rangle\).
Con~\eqref{eq:parabola},
\begin{align}
\langle u^2\rangle
&=
\frac{2}{R^2}\int_0^R
4\bar u^2\left(1-\frac{r^2}{R^2}\right)^2 r\,\dd r
=
8\bar u^2\int_0^1(1-\eta^2)^2\eta\,\dd\eta
\nonumber\\
&=
8\bar u^2\int_0^1\bigl(1-2\eta^2+\eta^4\bigr)\eta\,\dd\eta
=
8\bar u^2\int_0^1\bigl(\eta-2\eta^3+\eta^5\bigr)\,\dd\eta
\nonumber\\
&=
8\bar u^2\left[\frac{\eta^2}{2}-\frac{2\eta^4}{4}+\frac{\eta^6}{6}\right]_0^1
=
8\bar u^2\left(\frac12-\frac12+\frac16\right)
=
\frac{4}{3}\bar u^2.
\label{eq:u2}
\end{align}
El factor \(4/3\) es el exceso de \(\langle u^2\rangle\) sobre \(\bar u^2\):
en la parábola el centro va al doble de la media, y el promedio del cuadrado
pesa ese centro.

\section{Esfuerzo axial de esta compresión}

Con \(K\) constante,~\eqref{eq:duP} da
\(\p u/\p z=-(u/K)\,\dd P/\dd z\).
Para el perfil~\eqref{eq:parabola} y \(\mu\) constante,
\begin{equation}
\frac{\p}{\p z}\!\left(\mu\frac{\p u}{\p z}\right)
=
\frac{\mu u}{K^2}\left(\frac{\dd P}{\dd z}\right)^2
-
\frac{\mu u}{K}\frac{\dd^2 P}{\dd z^2}.
\label{eq:axial}
\end{equation}
Frente al viscoseo radial~\eqref{eq:F}, el primer término de~\eqref{eq:axial}
está en la razón
\begin{equation}
\frac{\mu\bar u\,|\dd P/\dd z|^2/K^2}{8\mu\bar u/R^2}
=
\frac{R^2}{8}\,\frac{|\dd P/\dd z|^2}{K^2}.
\label{eq:razon}
\end{equation}
Con \(|\dd P/\dd z|\sim 3\times 10^4\,\mathrm{Pa\,m^{-1}}\),
\(K=1.5\times 10^{10}\,\mathrm{Pa}\) y \(R=16\,\mathrm{m}\),
\eqref{eq:razon} es del orden de \(10^{-10}\).
El término con \(\dd^2P/\dd z^2\) queda en ese orden mientras la pendiente
varíe en la escala de la columna. Se retira~\eqref{eq:axial} de~\eqref{eq:prom}.

\section{La pendiente}

\begin{proposicion}
Bajo~\eqref{eq:phi0},~\eqref{eq:bc} y la hipótesis~\ref{hip:cierre},
la pendiente de presión del tramo sin burbujas es
\begin{equation}
\frac{\dd P}{\dd z}
=
-\,
\frac{\rho_m g+8\mu\bar u/R^2}
{1-\dfrac{4}{3}\,\bar u^2/c^2},
\qquad
c^2=\frac{K}{\rho_m},
\qquad
\bar u=\frac{\bar q}{\rho_m}.
\label{eq:dpdz}
\end{equation}
\end{proposicion}

\begin{proof}
Se sustituyen~\eqref{eq:F} y~\eqref{eq:u2} en~\eqref{eq:prom},
retirado el esfuerzo axial~\eqref{eq:axial}:
\begin{equation}
\frac{\dd P}{\dd z}
\left(1-\frac{4}{3}\frac{\bar u^2}{c^2}\right)
=
-\rho_m g-\frac{8\mu\bar u}{R^2}.
\end{equation}
Despejar \(\dd P/\dd z\) da~\eqref{eq:dpdz}.
La relación \(\bar u=\bar q/\rho_m\) es~\eqref{eq:ubar}.
\end{proof}

El numerador es el peso más el roce de pared de la parábola.
El denominador es la inercia de la compresión, después de promediar
\(u(r)^2\). Si en~\eqref{eq:u2} se reemplaza el perfil por una velocidad
plana, \(\langle u^2\rangle=\bar u^2\) y el denominador queda
\(1-\bar u^2/c^2\), que es el factor escrito en la recta algebraica del 1D.

\section{La recta de coeficientes congelados}

En el tramo algebraico del 1D se congelan \(\rho_m\), \(\mu\) y
\(\bar u=v_{\mathrm{in}}\), con \(K=15\,\mathrm{GPa}\).
Entonces~\eqref{eq:dpdz} es constante y
\begin{equation}
P(z)=P(H)+\frac{\dd P}{\dd z}\,(z-H),
\label{eq:recta}
\end{equation}
hasta la altura definida por
\begin{equation}
P(z_{\sat})
=
P_{\sat}
=
\left(\frac{c_0}{C_1}\right)^{1/\beta}.
\label{eq:Psat}
\end{equation}
Esa constante es la línea \texttt{dpdzcalc} del código 1D.
El denominador del código es \(1-v_{\mathrm{in}}^2/c^2\).
El de~\eqref{eq:dpdz} es \(1-(4/3)\bar u^2/c^2\), porque el promedio se hace
sobre la parábola.

\begin{obs}
En la base de Calbuco, con \(\bar u=18.3\,\mathrm{m\,s^{-1}}\),
\(c=2468\,\mathrm{m\,s^{-1}}\) y \(\rho_m=2463\,\mathrm{kg\,m^{-3}}\),
se tiene \(\bar u^2/c^2=5.5\times 10^{-5}\).
El denominador de~\eqref{eq:dpdz} vale \(1-7.3\times 10^{-5}\).
Sobre una pendiente de \(2.677\times 10^4\,\mathrm{Pa\,m^{-1}}\),
el divisor corrige unos \(2\,\mathrm{Pa\,m^{-1}}\).
\end{obs}

\begin{obs}[Dónde vale la parábola]
El dato de la base puede ser un tapón, \(u=v_{\mathrm{in}}\) para \(r<R\)
y \(u(R)=0\). Ese tapón cumple~\eqref{eq:bc} y tiene laplaciano nulo en el
interior, así que~\eqref{eq:parabola} lo describe una vez que el viscoseo
radial ya acomodó el perfil. La longitud de ese acomodo es el balance entre
inercia y viscoseo radial, \(L\sim\rho_m v_{\mathrm{in}} R^2/\mu\).
Con los números de la base queda del orden de \(2.5\,\mathrm{km}\).
La fórmula~\eqref{eq:dpdz} describe el tramo sin burbujas ya desarrollado,
desde ahí hasta \(z_{\sat}\). En Calbuco \(z_{\sat}\) cae cerca de
\(-3.7\,\mathrm{km}\).
\end{obs}

\section{El metro donde empieza la exsolución}
\label{sec:eps}

En \(z_{\sat}\) se tiene \(n=0\) y \(\varphi=0\). La velocidad del gas,
escrita como caudal de gas sobre \(\varphi\rho_g\), queda indeterminada.
El estado inicial de la marcha diferencial se arma un metro más arriba.
Se toma \(\varepsilon=1\,\mathrm{m}\). Si la saturación cae a menos de un
metro de la boca, \(\varepsilon=-z_{\sat}/10\). El punto adicional es
\begin{equation}
z_{\ad}=z_{\sat}+\varepsilon,
\qquad
P_{\ad}
=
P_{\sat}+\frac{\dd P}{\dd z}\,\varepsilon.
\label{eq:Pad}
\end{equation}
Como \(\dd P/\dd z<0\), se tiene \(P_{\ad}<P_{\sat}\) y Henry ya suelta agua.
La fracción másica exsuelta, con \(\xi\) todavía igual a \(\xi_0\), es
\begin{equation}
n_{\ad}
=
(1-\xi_0)\,
\frac{c_0-C_1 P_{\ad}^{\beta}}{1-C_1 P_{\ad}^{\beta}}.
\label{eq:n}
\end{equation}
\(P_{\ad}\) y \(\xi_0\) no dependen de \(r\), así que \(n_{\ad}\) es
constante en la sección. La densidad del gas es la del gas ideal,
\(\rho_{g,\ad}=P_{\ad}/(\Rv T)\).

\begin{proposicion}
\label{prop:ad}
En \(z_{\ad}\), con el caudal \(q(r)\) heredado del tramo sin burbujas,
\begin{align}
\varphi_{\ad}
&=
\left(
1+\frac{P_{\ad}}{n_{\ad}\Rv T}\,\frac{1-n_{\ad}}{\rho_m}
\right)^{-1},
\label{eq:phi}\\[0.4em]
u_m(r)
&=
u_g(r)
=
\frac{q(r)\,(1-n_{\ad})}{(1-\varphi_{\ad})\,\rho_m}
=
2\bar u_{\ad}\left(1-\frac{r^2}{R^2}\right),
\label{eq:velocidades}
\end{align}
donde
\begin{equation}
\bar u_{\ad}
=
\frac{\bar q\,(1-n_{\ad})}{(1-\varphi_{\ad})\,\rho_m}
=
\frac{\bar q\, n_{\ad}}{\varphi_{\ad}\,\rho_{g,\ad}},
\qquad
q(r)=2\bar q\left(1-\frac{r^2}{R^2}\right).
\label{eq:uad}
\end{equation}
En la pared, \(q(R)=0\), así que \(u_m(R)=u_g(R)=0\).
\(N\) y \(\xi\) siguen en los valores de la base, \(N_0\) y \(\xi_0\).
\end{proposicion}

\begin{proof}
La suma de las dos masas, con \(u_r=0\), conserva el caudal de cada radio.
El valor que trae el tramo sin burbujas es \(q(r)\) de~\eqref{eq:q}:
\begin{equation}
\rho_m(1-\varphi)\,u_m+\rho_g\,\varphi\,u_g=q(r).
\label{eq:qtotal}
\end{equation}
En este punto se parte el caudal con la fracción de Henry~\eqref{eq:n},
\begin{equation}
\rho_g\,\varphi\,u_g=n_{\ad}\,q(r),
\qquad
\rho_m(1-\varphi)\,u_m=(1-n_{\ad})\,q(r).
\label{eq:reparto}
\end{equation}
El volumen de gas por unidad de masa es \(n_{\ad}/\rho_{g,\ad}\) y el del
fundido es \((1-n_{\ad})/\rho_m\). Su cociente es \(\varphi_{\ad}\), que es
\eqref{eq:phi} al escribir \(\rho_{g,\ad}=P_{\ad}/(\Rv T)\).
De esa definición,
\begin{equation}
\frac{1-\varphi_{\ad}}{\varphi_{\ad}}
=
\frac{\rho_{g,\ad}\,(1-n_{\ad})}{n_{\ad}\,\rho_m},
\end{equation}
y por tanto
\begin{equation}
\frac{1-n_{\ad}}{(1-\varphi_{\ad})\,\rho_m}
=
\frac{n_{\ad}}{\varphi_{\ad}\,\rho_{g,\ad}}.
\label{eq:igualdad}
\end{equation}
Las dos velocidades de~\eqref{eq:reparto} coinciden:
\begin{equation}
u_m(r)
=
\frac{q(r)\,(1-n_{\ad})}{(1-\varphi_{\ad})\,\rho_m}
=
\frac{q(r)\,n_{\ad}}{\varphi_{\ad}\,\rho_{g,\ad}}
=
u_g(r).
\end{equation}
Con la parábola \(q(r)=2\bar q\,(1-r^2/R^2)\) se obtiene~\eqref{eq:velocidades}
y~\eqref{eq:uad}. El nodo \(r=R\) hereda \(q(R)=0\).
\end{proof}

\begin{obs}
\(\varphi_{\ad}\) no depende de \(r\). Las dos velocidades sí: conservan la
forma del caudal y se anulan en la pared. El deslizamiento \(u_g-u_m\)
empieza recién al integrar el momento desde \(z_{\ad}\); en este punto la
igualdad~\eqref{eq:igualdad} lo deja en cero.
\end{obs}

\end{document}
